\chapter{Introducción}

La descarga electrostática, conocida comúnmente como ESD por sus siglas en inglés
(\textit{Electrostatic Discharge}), representa un fenómeno crítico en ambientes donde se
manipulan, validan o caracterizan dispositivos electrónicos sensibles \cite{ESDAPrinciples2}.
En la industria de semiconductores y en laboratorios de electrónica, una descarga
electrostática puede provocar fallas inmediatas, degradación paramétrica o daños latentes
que afectan la confiabilidad de los dispositivos. Este problema adquiere mayor relevancia
debido a que muchas descargas no son perceptibles para el ser humano, pero pueden ser
suficientes para dañar circuitos integrados modernos \cite{TI2014}.

En la industria de semiconductores, la relevancia de este fenómeno aumenta conforme los
circuitos integrados reducen sus dimensiones y operan a mayores velocidades. Esta
tendencia incrementa la susceptibilidad de los dispositivos frente a eventos de descarga
electrostática, por lo que el control ESD se considera un aspecto crítico durante el diseño,
fabricación, manipulación, prueba y ensamblaje de circuitos integrados \cite{ESDA1}.
Desde el punto de vista de confiabilidad, una descarga electrostática puede causar fallas
inmediatas, pero también daños latentes que no necesariamente se detectan durante una
prueba inicial. Estos daños pueden degradar progresivamente el dispositivo hasta producir
una falla posterior, lo cual dificulta su detección y aumenta el riesgo de que el problema
aparezca durante etapas avanzadas de validación o incluso en campo \cite{TI2014}.

El impacto económico de los eventos ESD no se limita al costo directo del componente
dañado. También puede involucrar costos asociados a reprocesos, análisis de falla,
reparación, retrabajo, tiempo de ingeniería, pérdida de rendimiento de producción
(\textit{yield}), garantías y afectación de la confiabilidad del producto final
\cite{versalogic_esd}. La ESD Association indica que los eventos ESD continúan afectando
el rendimiento de producción, el costo de manufactura, la calidad del producto, la
confiabilidad y la rentabilidad de las empresas del sector electrónico \cite{ESDA1}. Según
Versalogic Corporation, las empresas de alta tecnología pueden perder entre un 4 \% y un
6 \% de sus ganancias brutas anuales debido a problemas asociados con la electrostática
\cite{versalogic_esd}.

En laboratorios de electrónica, el impacto del ESD también es relevante. Aunque estos
entornos no siempre corresponden a líneas de producción de alto volumen, en ellos se
manipulan equipos costosos, tarjetas electrónicas sensibles, circuitos integrados,
prototipos y muestras de ingeniería. Una descarga electrostática puede causar daños
inmediatos o degradar el desempeño del dispositivo, lo que puede generar actividades
adicionales de diagnóstico, reemplazo y repetición de pruebas
\cite{ESDA1,TI2014}. Por esta razón, la aplicación de medidas de control ESD
contribuye a proteger los dispositivos y a mantener la continuidad de las actividades
del laboratorio \cite{ESDAPrinciples2}.

Para reducir este riesgo, normativas como ANSI/ESD S20.20 e IEC 61340-5-1 establecen
lineamientos para la implementación de programas de control ESD en áreas protegidas
contra descargas electrostáticas. Estas prácticas incluyen el uso de superficies disipativas,
control de materiales, puesta a tierra del personal, uso de pulseras antiestáticas y
verificación periódica de los elementos de protección utilizados por los operadores
\cite{ESDAPrinciples2}.

Dentro de estos mecanismos, la pulsera antiestática permite establecer una trayectoria
controlada para la disipación de carga desde el usuario hacia tierra. Sin embargo, la sola
disponibilidad de una pulsera no garantiza que el usuario la utilice correctamente ni que
la conexión eléctrica sea adecuada. Por esta razón, se emplean probadores de pulseras
antiestáticas para verificar si la resistencia eléctrica del conjunto usuario-pulsera se
encuentra dentro de un rango aceptable. Equipos básicos, como el Desco 19240, permiten
realizar esta verificación mediante una indicación simple de aprobación o rechazo, pero no
incorporan identificación de usuario, almacenamiento histórico de pruebas ni generación
automática de reportes \cite{Desco_19240}.

\section{Escenario de aplicación: laboratorio TEM de Intel Costa Rica}

El presente proyecto surge a partir de una oportunidad de mejora identificada en el
laboratorio TEM de Intel Costa Rica, un entorno donde se utilizan estaciones de trabajo
para realizar pruebas funcionales y de caracterización a circuitos integrados. En este tipo
de laboratorio, la protección contra descargas electrostáticas es especialmente importante
debido al costo de los equipos utilizados, la sensibilidad de los dispositivos manipulados y
la necesidad de mantener buenas prácticas de control ESD durante las actividades de
laboratorio.

Actualmente, el laboratorio cuenta con un probador comercial de pulseras antiestáticas y
pulseras disponibles para los usuarios. El procedimiento esperado consiste en que cada
persona que vaya a utilizar una estación de trabajo tome una pulsera, realice la prueba
correspondiente y continúe únicamente si el resultado es satisfactorio. No obstante,
cuando el proceso depende únicamente de la disciplina individual del usuario, resulta
difícil verificar si todos los usuarios ejecutan la prueba de manera sistemática antes de
utilizar las estaciones de trabajo.

Esta situación evidencia que el problema no se limita a disponer de un equipo de
verificación ESD, sino a contar con un mecanismo que permita dar seguimiento al
cumplimiento del procedimiento. Si una prueba no se asocia con un usuario, fecha, hora y
resultado, el laboratorio no cuenta con información suficiente para determinar quiénes
realizaron la prueba, con qué frecuencia se ejecutó el procedimiento o si existe una brecha
entre el uso real del laboratorio y los registros de verificación ESD.

Por lo tanto, el laboratorio TEM de Intel Costa Rica representa un escenario adecuado para
validar una estación de control ESD. En este contexto, el prototipo propuesto permitiría
complementar el procedimiento actual, basado en el uso de un probador de pulseras,
mediante identificación de usuario, almacenamiento de registros, generación de reportes y
una interfaz gráfica que guíe el flujo de prueba.

En este contexto, el impacto esperado del prototipo consiste en reducir la
dependencia exclusiva de la disciplina individual del usuario mediante un proceso
registrable y verificable. Al asociar cada prueba con un usuario, fecha, hora y
resultado, además de incluir las condiciones ambientales como información
complementaria, el sistema permitiría que el encargado del
laboratorio consulte reportes históricos y audite el cumplimiento del
procedimiento ESD. De esta forma, aunque el prototipo no incorpora control físico
de acceso, sí convierte la verificación de pulseras antiestáticas en un proceso
documentado y sujeto a revisión.

\section{Motivación de una solución de bajo costo}

Existen sistemas comerciales avanzados capaces de integrar funciones como verificación
ESD, identificación de usuario, registro de pruebas, control de acceso, comunicación en
red y generación de reportes. Entre estas soluciones se encuentran plataformas como
SmartLog Pro 2, PGT130.DT y la estación de control ESDMAN
\cite{SmartLogPro2} \cite{PGT130DT} \cite{ESDMAN_0016809}. Sin embargo, estos sistemas
suelen presentar costos elevados, superiores a los 2000 USD en algunas configuraciones
\cite{ESDMAN_0016809}, y una capacidad limitada de personalización cuando se comparan
con soluciones basadas en módulos electrónicos de bajo costo y arquitecturas abiertas.

En este sentido, el laboratorio TEM de Intel Costa Rica se plantea como un entorno real
para evaluar la utilidad del prototipo y, al mismo tiempo, como una oportunidad para
mejorar el seguimiento del procedimiento ESD actualmente disponible. La solución
propuesta busca transformar una prueba manual y aislada en un proceso registrable,
verificable y trazable, manteniendo una arquitectura accesible, modular y personalizable.

El sistema propuesto integra un probador comercial de pulseras antiestáticas como
subsistema principal de verificación ESD, identificación de usuarios mediante lectura de
código de barras y reconocimiento facial, interfaz gráfica táctil, almacenamiento local de
registros y generación de reportes históricos. También incorpora temperatura y humedad
relativa como información complementaria de cada prueba.
De esta manera, el proyecto busca demostrar que es posible desarrollar una estación
inteligente de control ESD de bajo costo que mejore la trazabilidad de las pruebas en
entornos de laboratorio.

\section{Planteamiento del problema}

Aunque existen sistemas comerciales para el control ESD en laboratorios de electrónica y
manufactura de semiconductores, estas soluciones suelen presentar costos elevados y una
capacidad limitada de personalización. Adicionalmente, no se identifican alternativas de
bajo costo que integren de manera conjunta la verificación de pulseras antiestáticas, la
identificación de usuarios, el almacenamiento de registros históricos, la generación de
reportes y una interfaz gráfica para interactuar con el sistema.

Esta situación evidencia la necesidad de desarrollar una solución accesible y adaptable que
permita mejorar la trazabilidad de las pruebas ESD en entornos de laboratorio, como el
laboratorio TEM de Intel Costa Rica, donde se ha identificado una oportunidad de mejora
en el seguimiento del cumplimiento de los procedimientos ESD.

\section{Estructura del documento}

El presente documento se encuentra organizado en ocho capítulos que describen de forma
progresiva la fundamentación, el diseño, la formulación, la estrategia de validación y la
planificación del proyecto.

En el Capítulo 1 se presentan los antecedentes relacionados con el control de descargas
electrostáticas, la motivación del proyecto, el escenario de aplicación en el laboratorio
TEM de Intel Costa Rica, la necesidad de una solución de bajo costo y el planteamiento
del problema.

El Capítulo 2 desarrolla el marco teórico necesario para comprender los conceptos
fundamentales asociados a las descargas electrostáticas, la medición de resistencia
cuerpo-tierra y de alta impedancia, el protocolo de comunicación ESP-NOW y los algoritmos
ligeros de visión artificial para el reconocimiento facial embebido.

En el Capítulo 3 se presenta el estado del arte, incluyendo investigaciones relacionadas,
normativas aplicables y un análisis comparativo de soluciones comerciales utilizadas para
el control ESD, identificando sus principales características y limitaciones.

El Capítulo 4 describe la solución propuesta, detallando la arquitectura general del
sistema, los subsistemas que la conforman, los dispositivos seleccionados, los requisitos
del sistema y la integración electrónica, funcional y física del prototipo.

El Capítulo 5 presenta el presupuesto estimado para la implementación del sistema y una
comparación económica con soluciones comerciales disponibles.

En el Capítulo 6 se establecen la hipótesis y los objetivos del proyecto, así como las
actividades, entregables, alcances y limitaciones considerados para su desarrollo.

El Capítulo 7 presenta la estrategia de validación del prototipo, incluyendo las
condiciones del experimento, los indicadores de validación, el procedimiento experimental
y los criterios de aceptación. La validación se planifica preferentemente en el laboratorio
TEM de Intel Costa Rica.

Finalmente, el Capítulo 8 presenta el cronograma de trabajo propuesto para la ejecución
del proyecto.
